\section{问题描述}

\subsection{AUV动力学模型}

AUV的6自由度动力学方程采用Fossen矢量表示\cite{fossen2011handbook}，基于相对速度形式：

\begin{equation}
    \bm{M}\dot{\bm{\nu}}_r + \bm{C}(\bm{\nu}_r)\bm{\nu}_r + \bm{D}(\bm{\nu}_r)\bm{\nu}_r + \bm{g}(\bm{\eta}) = \bm{\tau} + \bm{d}(t)
    \label{eq:dynamics}
\end{equation}

其中$\bm{\nu}_r = \bm{\nu} - \bm{\nu}_c$为相对水流速度，$\bm{\nu}$为体坐标系对地速度，$\bm{\nu}_c$为体坐标系海流分量，$\bm{M}$为惯性矩阵（刚体+附加质量），$\bm{C}(\bm{\nu}_r)$为科氏力矩阵，$\bm{D}(\bm{\nu}_r)$为水动力阻尼矩阵，$\bm{g}(\bm{\eta})$为重力/浮力向量，$\bm{\tau}$为控制力/力矩，$\bm{d}(t)$为未建模扰动。

水动力阻尼采用CFD RANS标定的二次型模型：

\begin{equation}
    \tau_{\text{drag},i} = -(d_{i1}\,\nu_{r,i} + d_{i2}\,\nu_{r,i}\,|\nu_{r,i}|), \quad i=1,\dots,6
    \label{eq:drag}
\end{equation}

其中$d_{i1}$、$d_{i2}$为CFD标定的线性和二次阻尼系数（表\ref{tab:cfd}）。

\begin{table}[h]
\centering
\caption{XHY AUV CFD标定阻力系数}
\label{tab:cfd}
\begin{tabular}{lcccccc}
\toprule
系数 & 纵荡 & 横荡 & 沉浮 & 横滚 & 俯仰 & 偏航 \
\hline
$d_1$ & 0.65 & 1.67 & 2.95 & --- & 0.5941 & 0.2046 \
$d_2$ & 10.85 & 178.59 & 40.83 & --- & 4.088 & 18.200 \
$R^2$ & $>$0.99 & $>$0.99 & $>$0.99 & --- & $>$0.97 & $>$0.98 \
\bottomrule
\end{tabular}
\end{table}

\subsection{海流参数化}

在NED惯性坐标系中，水平海流由速度幅值$V_c$和方向角$\beta_c$参数化：

\begin{equation}
    \bm{c} = \begin{bmatrix} c_N \ c_E \end{bmatrix} = \begin{bmatrix} V_c\cos\beta_c \ V_c\sin\beta_c \end{bmatrix}
    \label{eq:current_ned}
\end{equation}

在体坐标系下，海流分量随AUV姿态变化：

\begin{equation}
    \bm{\nu}_c = \begin{bmatrix} V_c\cos(\beta_c-\psi) \ V_c\sin(\beta_c-\psi) \ w_c \ 0 \ 0 \ 0 \end{bmatrix}
    \label{eq:current_body}
\end{equation}

其中$\psi$为航向角。海流的时间演化建模为一阶Gauss-Markov过程，相关时间常数$\tau_c=100$ s。

\subsection{CFD模型不确定性的量化表征}

CFD标定提供了名义阻力系数$\bar{d}_{ij}$，但真实阻力系数存在不确定性。对于给定的扰动水平$\delta$（$0\le\delta\le 0.5$），实际阻力系数为：

\begin{equation}
    d_{ij}^{\text{plant}} = \bar{d}_{ij} \cdot (1 + \delta \cdot U[-1,1])
    \label{eq:mismatch}
\end{equation}

其中$U[-1,1]$为均匀分布随机变量（固定种子可复现）。观测器始终使用名义系数$\bar{d}_{ij}$，plant使用扰动系数$d_{ij}^{\text{plant}}$。
